%============================================================
% CH1 · LESSON 3 — THE COST MODEL
%============================================================
\subsection{The Cost Model: Reading the Constraints Before Coding}
\label{ssec:ch1-cost}

A problem statement tells you far more than the task: its \textbf{constraints} --- the limits on $n$, on the values, on the number of test cases --- tell you how fast the algorithm must be and which integer type it needs. Reading them \emph{before} you think about the algorithm saves you from two classic failures: a correct solution that gets \tle{}, and a correct solution that gets \wa{} because a sum silently overflowed.

%------------------------------------------------------------
\subsubsection*{1. The budget: about $10^8$ simple operations per second}

A judge executes roughly $10^8$ ``simple'' operations per second --- an addition, a comparison, an array access with a sequential pattern (Lesson~\ref{ssec:ch1-array}). With a typical limit of 1--2 seconds, your algorithm may perform a few times $10^8$ steps. This is a rule of thumb, not a law: a loop doing only additions over an array may run $10^9$ steps per second, while a loop with divisions, \texttt{map} lookups or random memory accesses may manage only $10^7$.

To use the budget, count the steps of your algorithm as a function of $n$ and plug in the largest $n$:

\begin{table}[H]
\centering\small
\renewcommand{\arraystretch}{1.2}
\begin{tabular}{@{}lllll@{}}
\toprule
\textbf{loop shape} & \textbf{steps} & $n = 10^3$ & $n = 10^5$ & $n = 2\cdot10^5$ \\
\midrule
one scan & $n$ & $10^3$ & $10^5$ & $2\cdot10^5$ \\
sort, then scan & $n \log_2 n$ & $10^4$ & $1.7\cdot10^6$ & $3.5\cdot10^6$ \\
all pairs $i < j$ & $n(n-1)/2$ & $5\cdot10^5$ & $5\cdot10^9$ \;\tle & $2\cdot10^{10}$ \;\tle \\
all subarrays with an inner loop & $\approx n^3/6$ & $1.7\cdot10^8$ & --- \;\tle & --- \;\tle \\
\bottomrule
\end{tabular}
\end{table}

The front-matter chapter \emph{From Constraints to Algorithm} (Chapter~\ref{ch:constraints}) turns this into a lookup table. The three rows you need constantly:

\begin{diagrambox}
\begin{tikzpicture}[x=1cm, y=1cm,
    lim/.style={rectangle, rounded corners=3pt, draw=primaryblue, fill=theorybg, line width=0.9pt,
                font=\scriptsize, align=center, text width=2.7cm, minimum height=1.25cm},
    alg/.style={rectangle, rounded corners=3pt, draw=successgreen, fill=solutionbg, line width=0.9pt,
                font=\scriptsize, align=center, text width=3.8cm, minimum height=1.25cm}]
    \node[lim] (a) at (0,0)    {$n \le 5000$};
    \node[lim] (b) at (0,-1.6) {$n \le 2\cdot10^5$};
    \node[lim] (c) at (0,-3.2) {$n \le 10^6$--$10^7$};
    \node[alg] (A) at (4.4,0)    {$\bigO{n^2}$: all pairs, every subarray with running state};
    \node[alg] (B) at (4.4,-1.6) {$\bigO{n \log n}$: sort, binary search; or $\bigO{n}$};
    \node[alg] (C) at (4.4,-3.2) {$\bigO{n}$ only, with fast input};
    \foreach \x/\y in {a/A, b/B, c/C} \draw[flow] (\x) -- (\y);
    \node[note, anchor=west, text width=3.0cm, align=left] at (7.1,-1.6)
        {This chapter's problems have $n \le 10^4$--$2\cdot10^5$: every solution here is a single scan or two, $\bigO{n}$.};
\end{tikzpicture}
\figcap{From the limit on $n$ to the complexity class. Read $n$ first; it rules out most ideas immediately.}
\end{diagrambox}

\begin{examplebox}[Reading a statement]
``$1 \le t \le 10^4$ test cases, $2 \le n \le 2\cdot10^5$, the sum of $n$ over all test cases does not exceed $2\cdot10^5$, $1 \le a_i \le 10^9$.''\\[2pt]
$\Rightarrow$ total input size $2\cdot10^5$: $\bigO{n \log n}$ or $\bigO{n}$ per test; $\bigO{n^2}$ is $4\cdot10^{10}$, far too slow. Values up to $10^9$ fit in \texttt{int}, but a sum of up to $2\cdot10^5$ of them reaches $2\cdot10^{14}$: \texttt{long long}. Many test cases: read fast and do not reset big arrays per test (Section~3).
\end{examplebox}

%------------------------------------------------------------
\subsubsection*{2. Overflow: \texttt{int} versus \texttt{long long}}

\begin{table}[H]
\centering\small
\renewcommand{\arraystretch}{1.2}
\begin{tabularx}{\linewidth}{@{}l l >{\raggedright\arraybackslash}X@{}}
\toprule
\textbf{type} & \textbf{range} & \textbf{rule of thumb} \\
\midrule
\texttt{int} (32 bits) & $-2^{31} \dots 2^{31}-1 = 2147483647$ & $\approx 2.1\cdot10^9$: values up to $10^9$ fit, their sums do not \\
\texttt{long long} (64 bits) & $-2^{63} \dots 2^{63}-1 = 9223372036854775807$ & $\approx 9.2\cdot10^{18}$: sums and products of two values up to $10^9$ \\
\texttt{unsigned long long} & $0 \dots 2^{64}-1$ & wraps modulo $2^{64}$ (hashing) \\
\bottomrule
\end{tabularx}
\end{table}

Signed overflow in C++ is \emph{undefined behaviour}: the program may print garbage, loop forever, or appear to work on your machine and fail on the judge. In practice an \texttt{int} sum wraps around modulo $2^{32}$. The program below shows what the wrapped results look like (it computes them with unsigned arithmetic, which is well defined):

\begin{lstlisting}[style=cpp]
#include <bits/stdc++.h>
using namespace std;

int main() {
    cout << "INT_MAX   = " << INT_MAX << "\n";
    cout << "LLONG_MAX = " << LLONG_MAX << "\n";
    int n = 200000, v = 1000000000;            // n values, each 10^9
    int sumInt = 0;
    long long sumLL = 0;
    for (int i = 0; i < n; i++) {
        sumInt = (int)((unsigned)sumInt + (unsigned)v);   // what an int sum wraps to
        sumLL += v;
    }
    cout << "int sum       = " << sumInt << "\n";         // 552894464  (garbage)
    cout << "long long sum = " << sumLL << "\n";          // 200000000000000
    int a = 100000, b = 100000;
    long long wrong = (long long)(int)((unsigned)a * (unsigned)b);  // what a*b wraps to in int
    long long right = (long long)a * b;                             // widen BEFORE multiplying
    cout << "a*b in int  = " << wrong << "\n";                      // 1410065408
    cout << "a*b in ll   = " << right << "\n";                      // 10000000000
    int lo = 2000000000, hi = 2100000000;
    cout << "lo + (hi - lo) / 2 = " << lo + (hi - lo) / 2 << "\n";  // 2050000000
    cout << "pairs for n = 2e5: " << (long long)n * (n - 1) / 2 << "\n";   // 19999900000
}
\end{lstlisting}

\begin{diagrambox}
\begin{tikzpicture}[x=1cm, y=1cm]
    \draw[thick, darkgray] (0,0) -- (12,0);
    \foreach \x/\l in {0/{$-2^{31}$}, 6/0, 12/{$2^{31}-1$}}{ \draw[thick] (\x,0.12) -- (\x,-0.12); \node[font=\scriptsize, below] at (\x,-0.12) {\l}; }
    \fill[successgreen!25] (6,0.18) rectangle (8.8,0.42);
    \node[font=\scriptsize, text=successgreen!60!black] at (7.4,0.65) {values $\le 10^9$};
    \draw[warningred, line width=1.2pt, -{Stealth[length=2.5mm]}] (8.8,0.3) .. controls (11,1.4) and (12.6,1.2) .. (12.3,0.3);
    \draw[warningred, line width=1.2pt, dashed, -{Stealth[length=2.5mm]}] (-0.3,0.3) .. controls (-0.6,1.2) and (0.4,1.4) .. (1.6,0.35);
    \node[font=\scriptsize\bfseries, text=warningred] at (6,1.45) {past $2^{31}-1$ the sum re-enters at $-2^{31}$};
\end{tikzpicture}
\figcap{\texttt{int} arithmetic behaves like a clock with $2^{32}$ positions: one step past the largest value lands on the smallest.}
\end{diagrambox}

Three overflow traps cover almost every case:
\begin{enumerate}[leftmargin=*]
    \item \textbf{Sums.} $n$ values of size up to $V$ sum up to $nV$. With $n = 2\cdot10^5$ and $V = 10^9$ that is $2\cdot10^{14}$: the accumulator must be \texttt{long long}. The \texttt{int} sum above prints $552894464$ --- a plausible-looking, completely wrong number.
    \item \textbf{Products are computed in the type of the operands.} \texttt{long long x = a * b;} with \texttt{int a, b} multiplies in \texttt{int} \emph{first} and only then converts the garbage: $10^5 \cdot 10^5$ becomes $1410065408$. Write \texttt{(long long)a * b} or declare the values \texttt{long long}.
    \item \textbf{Counts of pairs and subarrays.} $n(n-1)/2$ for $n = 2\cdot10^5$ is $19999900000 \approx 2\cdot10^{10}$: \texttt{long long}, and compute it as \texttt{(long long)n * (n - 1) / 2}.
\end{enumerate}
A fourth, rarer one: the midpoint \texttt{(lo + hi) / 2} overflows when both are near $2^{31}$; for \texttt{lo} $= 2\cdot10^9$, \texttt{hi} $= 2.1\cdot10^9$ the \texttt{int} sum wraps to $-194967296$ and the ``midpoint'' becomes $-97483648$. \texttt{lo + (hi - lo) / 2} $= 2050000000$ is always safe.

\begin{insightbox}[When in doubt]
Use \texttt{long long} for every sum, product, count and answer. It costs nothing measurable in a scan; an overflow costs the whole problem. Keep \texttt{int} for indices and for large arrays whose values fit (memory: \texttt{long long} is 8 bytes).
\end{insightbox}

%------------------------------------------------------------
\subsubsection*{3. Many test cases: the ``sum of $n$'' guarantee}

Codeforces problems usually pack many test cases into one input: ``$1 \le t \le 10^4$, $1 \le n \le 2\cdot10^5$, \emph{the sum of $n$ over all test cases does not exceed $2\cdot10^5$}''. The last sentence is the important one. Without it, the input could hold $10^4$ tests of size $2\cdot10^5$ each ($2\cdot10^9$ numbers --- impossible to even read). With it, the \emph{total} work of an $\bigO{n}$ algorithm over all tests is $\bigO{\sum n} = \bigO{2\cdot10^5}$.

\begin{diagrambox}
\begin{tikzpicture}[x=1cm, y=1cm]
    \node[cell, fill=hitfill, minimum width=5mm] (t1) at (0.25,0) {};
    \node[cell, fill=basefill, minimum width=24mm, right=0pt of t1] (t2) {$n_2$};
    \node[cell, fill=hitfill, minimum width=4mm, right=0pt of t2] (t3) {};
    \node[cell, fill=basefill, minimum width=12mm, right=0pt of t3] (t4) {$n_4$};
    \node[cell, fill=hitfill, minimum width=18mm, right=0pt of t4] (t5) {$n_5$};
    \node[cell, fill=basefill, minimum width=3mm, right=0pt of t5] (t6) {};
    \draw[thick, primaryblue, decorate, decoration={brace, mirror, amplitude=5pt}]
        ([yshift=-2mm]t1.south west) -- ([yshift=-2mm]t6.south east)
        node[midway, below=5pt, font=\scriptsize\bfseries, text=primaryblue] {$n_1 + n_2 + \dots + n_t \le 2\cdot10^5$};
    \node[note, anchor=west, text width=5.2cm, align=left] at (7.8,0)
        {a few big tests or many tiny ones --- the total length is bounded, so $\bigO{n}$ per test is $\bigO{\sum n}$ overall};
\end{tikzpicture}
\figcap{The tests of one input laid end to end: their total length is what the limit bounds.}
\end{diagrambox}

The guarantee covers work proportional to $n$, \emph{not} work proportional to the maximal $n$ or to the value range:

\begin{warningbox}[Resetting a global array in every test]
\begin{lstlisting}[style=snippet]
int cnt[200005];
while (t--) {
    memset(cnt, 0, sizeof cnt);            // WRONG: 2e5 per test, 1e4 tests -> 2e9 writes
    ...
}
while (t--) {
    read n and a[0..n);
    for (int i = 0; i < n; i++) cnt[a[i]]++;
    ...
    for (int i = 0; i < n; i++) cnt[a[i]] = 0;   // RIGHT: undo only what this test touched
}
\end{lstlisting}
The same applies to \texttt{vector<int> cnt(MAXV)} declared inside the loop: allocate \texttt{cnt(n + 1)} sized by this test's $n$ (or by its values), or clear only the used cells.
\end{warningbox}

%------------------------------------------------------------
\subsubsection*{4. Fast input and output}

Reading $2\cdot10^5$ numbers with plain \texttt{cin} is fine; reading $10^6$ or more can take a noticeable part of the time limit, because by default \texttt{cin} synchronises with C's \texttt{stdio} and flushes \texttt{cout} before every read. Two lines at the top of \texttt{main} switch both off:

\begin{lstlisting}[style=cpp]
#include <bits/stdc++.h>
using namespace std;

int main() {
    ios::sync_with_stdio(false);   // stop syncing with printf/scanf (do not mix them now)
    cin.tie(nullptr);              // do not flush cout before every cin
    int t;
    if (!(cin >> t)) return 0;
    while (t--) {
        int n;
        cin >> n;
        vector<long long> a(n);
        long long sum = 0;
        for (auto& x : a) { cin >> x; sum += x; }
        cout << sum << '\n';       // '\n', not endl: endl flushes the buffer every time
    }
}
\end{lstlisting}

On input \texttt{2 / 3 / 1 2 3 / 2 / 1000000000 1000000000} it prints $6$ and $2000000000$ --- the second sum already exceeds \texttt{INT\_MAX}, which is why \texttt{a} holds \texttt{long long}.

\begin{warningbox}[\texttt{endl} in a loop]
\texttt{cout << x << endl} flushes the output buffer on every line. With $2\cdot10^5$ output lines this alone can cause \tle{}. Use \texttt{'\textbackslash n'} and let the buffer flush once at the end. (Interactive problems are the exception: there you \emph{must} flush after each query.)
\end{warningbox}

%------------------------------------------------------------
\subsubsection*{5. Memory}

A typical limit is 256 MB. An \texttt{int} is 4 bytes, so $10^7$ \texttt{int}s take 40 MB and $6\cdot10^7$ take 240 MB. A $5000 \times 5000$ \texttt{int} table is 100 MB (fine); a $10^4 \times 10^4$ one is 400 MB (\mle). \texttt{vector<vector<int>>} adds about 24 bytes of bookkeeping per row; \texttt{vector<bool>} packs 8 flags per byte. In this chapter memory is never tight, but it decides later whether an $\bigO{n^2}$ table is allowed even when the time is.

\begin{checklistbox}[Before you code: read the constraints]
\begin{itemize}[leftmargin=*]
    \item[$\square$] Largest $n$ (and $\sum n$): which complexity fits in $\approx 10^8$ steps per second? (Chapter~\ref{ch:constraints})
    \item[$\square$] Largest value $V$: does every value fit in \texttt{int}? Does every \emph{sum}, \emph{product}, \emph{count} and the \emph{answer}?
    \item[$\square$] Products: is one operand cast to \texttt{long long} \emph{before} the multiplication?
    \item[$\square$] Many tests: is per-test work proportional to this test's $n$ (no full \texttt{memset})?
    \item[$\square$] Large input: \texttt{sync\_with\_stdio(false)}, \texttt{cin.tie(nullptr)}, \texttt{'\textbackslash n'}.
    \item[$\square$] Memory: number of cells $\times$ bytes per cell below the limit?
\end{itemize}
\end{checklistbox}
